% Chapter 4: approach, data and evaluation for each research question.
\chapter{Methodology}\label{ch:methods}

\section{Temperature model}

Let $T(\mathbf{s},t)$ be air temperature at site $\mathbf{s}$ and hour
$t$. The model writes it as a covariate trend plus a spatially
correlated residual,
\begin{equation}
  T(\mathbf{s},t) = \mathbf{x}(\mathbf{s})^{\top}\boldsymbol{\beta}_t
  + f(\mathbf{s},t) + \varepsilon,
  \qquad f \sim \mathcal{GP}\bigl(0, k_{\theta}\bigr),
  \label{eq:model}
\end{equation}
where $\mathbf{x}$ holds sky-view factor, impervious fraction, canopy
cover and distance to water at 3\,m resolution, and $k_\theta$ is a
product of a Matérn kernel in space and a periodic kernel in hour of
day. The model is fitted with stochastic variational inference on 512
inducing points.

\section{Exposure-weighted acquisition (RQ1)}

For a candidate site $\mathbf{c}$, the acquisition score is the
reduction in exposure-weighted predictive variance across the grid
$\mathcal{G}$,
\begin{equation}
  a(\mathbf{c}) = \sum_{\mathbf{g}\in\mathcal{G}} w(\mathbf{g})
  \bigl[\sigma^2(\mathbf{g}) - \sigma^2(\mathbf{g}\mid\mathbf{c})\bigr],
  \label{eq:acq}
\end{equation}
with weights $w$ from the census vulnerability index (residents over
75, living alone, top-floor flats). Batches are chosen greedily, which
keeps the near-optimality guarantee of the unweighted case.

\textbf{Hypothesis H1.} At equal sensor counts, rule \eqref{eq:acq}
lowers exposure-weighted RMSE by at least 20\,\% relative to
variance-only placement, with no more than 5\,\% loss in unweighted RMSE.

\section{Transfer between cities (RQ2)}

The model trained on Aldmere is applied to two partner cities with
dense reference networks, Lyre and Westhaven, after withholding all
their data. Placement is simulated from zero to 150 sensors, comparing
a cold start, a transferred prior and a prior adapted by importance
weighting on covariate distributions \parencite{sugiyama2007covariate}.

\textbf{Hypothesis H2.} A transferred prior halves the number of
sensors needed to reach 1.0\,K weighted RMSE in the target city.

\section{Field validation (RQ3)}

Sixty battery-powered units with radiation shields are deployed from
June to September 2027. Every two weeks, the twenty with the lowest
marginal value are moved to the sites the rule proposes, subject to a
list of permitted mounting points agreed with the city.

\textbf{Hypothesis H3.} Field gains are at least two thirds of the gains
predicted by simulation for the same period.
